%% kapa: exam=1 round=20 year=2009 period=1 form=A subject=civil_law subject_ko=민법 seq=20 number=20
%% source: Q-Net 감정평가사 기출문제 (cst00309) artlSeq=5024178 file=extracted/1008/5024178/2009년 제20회 감정평가사 1차시험 기출문제/1교시.hwp
%% conversion: kordoc 4.24.51 -> tools/gen1.py (auto); review=proofread-sonnet; flags=none
\begin{question}{20}
선의취득에 관한 설명으로 옳지 \uline{않은} 것은? (다툼이 있으면 판례에 의함)
\begin{choices}
  \item 선의취득에 의한 물권취득의 효력은 법률의 규정에 의한 것이므로, 선의취득자가 임의로 선의취득의 효과를 거부하고 종전 소유자에게 목적물을 반환받아 갈 것을 요구할 수 없다.
  \item 타인의 동산을 보관하던 乙을 소유자라고 오신하여 그로부터 동산을 매수함과 동시에 그것을 乙에게 임대해 준 甲은 그 동산을 선의취득한다.
  \item 횡령된 동산도 선의취득의 대상이 될 수 있다.
  \item 채권자가 질권설정자에게 처분권이 있다고 믿고 과실 없이 질권을 설정받은 경우에는 질권을 선의취득한다.
  \item 채무자나 물상보증인의 소유가 아닌 동산을 경매절차에서 매각대금을 완납하고 인도받은 매수인(경락인)은 그 동산을 선의취득할 수 있다.
\end{choices}
\end{question}
\answer{2}
